\chapter{Inferential Statistics in Full}
\label{app:inferential}

This appendix reproduces the complete inferential test output
referenced in Chapter~\ref{ch:evaluation} and
Chapter~\ref{ch:discussion}. The within-subjects user-study tests
use $n = 69$ (matched within-subjects, Friedman + Wilcoxon); the
document-quality and timing analyses use $n = 113$ (between-groups,
Kruskal--Wallis + Mann--Whitney $U$); the confidence--quality
correlations supporting the RQ3 discussion use the joined
document-survey sample $n = 102$. The full sample-size derivation
is reproduced first, followed by the user-study inferential tests,
the NLP inferential tests, the confidence--quality correlations,
and the timing inferential tests.

\section{Sample-Size Derivation}
\label{app:samplesizes}

All sample sizes used in the thesis are derived from the raw data
files committed to the repository:
\texttt{python/survey/docs/data.csv} for the user study (one row per
participant) and \texttt{python/metrics/output/all\_scores\_combined.csv}
for the NLP and judge scores (one row per LLM-generated document).

\begin{table}[htbp]
\centering
\small
\begin{tabular}{lr}
\toprule
\textbf{User-study subset} & \textbf{n} \\
\midrule
Recruited participants (raw) & 73 \\
Participants with at least one valid session ($p\_0001 \geq 1$) & 69 \\
Form responses (mode-level) & 69 \\
Ask responses (mode-level) & 69 \\
Chat responses (mode-level) & 69 \\
Matched-triple participants (valid response in all three modes) & 69 \\
\bottomrule
\end{tabular}
\caption{User-study sample sizes after consent and quality filtering.}
\label{tab:n-survey}
\end{table}

\begin{table}[htbp]
\centering
\small
\begin{tabular}{lr}
\toprule
\textbf{NLP / judge subset} & \textbf{n} \\
\midrule
Total documents with parseable NLP scores and judge ratings & 113 \\
Form documents & 37 \\
Ask documents & 37 \\
Chat documents & 39 \\
Unique participants represented & 47 \\
Participants with at least one document in each mode & 28 \\
Strict matched triple (exactly one document per mode) & 26 \\
\bottomrule
\end{tabular}
\caption{NLP and judge sample sizes. The full $n = 113$ sample is
used for the document-quality inferential tests (Kruskal--Wallis,
Mann--Whitney $U$).}
\label{tab:n-nlp}
\end{table}

\begin{table}[htbp]
\centering
\small
\begin{tabular}{lr}
\toprule
\textbf{User-study $\times$ NLP join} & \textbf{n} \\
\midrule
Raw joined rows (document with matching user-study response) & 106 \\
Deduplicated joined rows (one row per participant--mode) & 102 \\
Form joined rows (deduplicated) & 34 \\
Ask joined rows (deduplicated) & 33 \\
Chat joined rows (deduplicated) & 35 \\
Unique participants in the joined sample & 44 \\
Participants with user-study $\times$ NLP coverage in all three modes & 27 \\
\bottomrule
\end{tabular}
\caption{Triangulated samples used for confidence--quality
analyses. The canonical RQ3 sample is the deduplicated join
$n = 102$. Four participants had multiple submissions per mode in
\texttt{document\_results.csv}; the canonical row per
participant--mode was selected by highest \texttt{id} (latest
submission). Three participants in the documents file had no
matching survey row and were dropped. The within-subjects
$n = 27$ is used for paired confidence~$\times$~quality
correlations.}
\label{tab:n-triangulation}
\end{table}

\section{User-Study Inferential Tests}
\label{app:inf-survey}

\subsection{Normality (Shapiro--Wilk, $n = 69$)}

\begin{table}[htbp]
\centering
\caption{Shapiro--Wilk normality tests for user-study measures by mode.}
\label{tab:app-survey-normality}
\small
\begin{tabular}{@{}l l r r l@{}}
\toprule
Measure & Mode & $W$ & $p$ & Normal? \\
\midrule
SUS & Form & 0.944 & 0.004 & No \\
SUS & Ask  & 0.976 & 0.210 & Yes \\
SUS & Chat & 0.927 & 0.001 & No \\
R-TLX composite & Form & 0.940 & 0.003 & No \\
R-TLX composite & Ask  & 0.926 & 0.001 & No \\
R-TLX composite & Chat & 0.945 & 0.005 & No \\
AI Support (6-item) & Form & 0.923 & $<0.001$ & No \\
AI Support (6-item) & Ask  & 0.965 & 0.048 & No \\
AI Support (6-item) & Chat & 0.952 & 0.010 & No \\
AI Confidence ($5_{ai}$) & Form & 0.846 & $<0.001$ & No \\
AI Confidence ($5_{ai}$) & Ask  & 0.856 & $<0.001$ & No \\
AI Confidence ($5_{ai}$) & Chat & 0.864 & $<0.001$ & No \\
AI Reuse ($6_{ai}$) & Form & 0.799 & $<0.001$ & No \\
AI Reuse ($6_{ai}$) & Ask  & 0.872 & $<0.001$ & No \\
AI Reuse ($6_{ai}$) & Chat & 0.861 & $<0.001$ & No \\
\bottomrule
\end{tabular}
\end{table}

Non-normal distributions are detected for nearly every measure. Non-parametric
tests (Friedman omnibus and Wilcoxon signed-rank post-hoc) are used
throughout.

\subsection{Friedman Omnibus ($n = 69$ matched triple)}

\begin{table}[htbp]
\centering
\caption{Friedman omnibus tests across the three modes for user-study measures.
Kendall's $W$: 0--0.1 weak, 0.1--0.3 moderate, 0.3--0.5 strong.}
\label{tab:app-survey-friedman}
\small
\begin{tabular}{@{}l r r r r l@{}}
\toprule
Measure & $\chi^2$ & df & $p$ & $W$ & Significant \\
\midrule
SUS (0--100)                    & 5.51  & 2 & 0.064  & 0.040 & No \\
R-TLX composite (1--21)         & 5.24  & 2 & 0.073  & 0.038 & No \\
R-TLX Mental Demand             & 3.13  & 2 & 0.209  & 0.023 & No \\
R-TLX Physical Demand           & 2.08  & 2 & 0.354  & 0.015 & No \\
R-TLX Temporal Demand           & 12.77 & 2 & 0.002  & 0.093 & \textbf{Yes} \\
R-TLX Effort                    & 3.78  & 2 & 0.151  & 0.027 & No \\
R-TLX Frustration               & 8.93  & 2 & 0.012  & 0.065 & \textbf{Yes} \\
R-TLX Performance (pos.)        & 2.33  & 2 & 0.311  & 0.017 & No \\
AI Confidence ($5_{ai}$)        & 3.76  & 2 & 0.152  & 0.027 & No \\
AI Support (6-item composite)   & 1.05  & 2 & 0.592  & 0.008 & No \\
AI Reuse intention ($6_{ai}$)   & 11.12 & 2 & 0.004  & 0.081 & \textbf{Yes} \\
\bottomrule
\end{tabular}
\end{table}

\subsection{Pairwise Wilcoxon Signed-Rank Post-hoc (Bonferroni $\alpha = 0.0167$)}

Post-hoc tests are reported for measures with a significant Friedman result
and, for completeness, for SUS and the R-TLX composite. Sign convention:
positive $r$ means the first-named mode ranks higher. Effect-size labels:
$|r| < 0.1$ negligible, $0.1$--$0.3$ small, $0.3$--$0.5$ medium, $\geq 0.5$ large.

\begin{table}[htbp]
\centering
\caption{Pairwise Wilcoxon signed-rank post-hoc tests for user-study measures.}
\label{tab:app-survey-posthoc}
\small
\setlength{\tabcolsep}{4pt}
\begin{tabular}{@{}l l r r r r l@{}}
\toprule
Measure & Pair & $n$ & $W$ & $p$ & $r$ & Sig (Bonf) \\
\midrule
SUS               & Form vs Ask  & 60 & 578.5 & 0.013   & $+0.368$ medium      & \textbf{Yes} \\
SUS               & Form vs Chat & 58 & 549.5 & 0.018   & $+0.358$ medium      & No \\
SUS               & Ask vs Chat  & 61 & 885.5 & 0.666   & $-0.063$ negligible  & No \\
R-TLX composite   & Form vs Ask  & 65 & 684.0 & 0.011   & $-0.362$ medium      & \textbf{Yes} \\
R-TLX composite   & Form vs Chat & 65 & 636.0 & 0.004   & $-0.407$ medium      & \textbf{Yes} \\
R-TLX composite   & Ask vs Chat  & 63 & 995.5 & 0.932   & $+0.012$ negligible  & No \\
Temporal Demand   & Form vs Ask  & 43 & 237.0 & 0.004   & $-0.499$ medium      & \textbf{Yes} \\
Temporal Demand   & Form vs Chat & 46 & 233.0 & $<0.001$& $-0.569$ large       & \textbf{Yes} \\
Temporal Demand   & Ask vs Chat  & 46 & 427.0 & 0.211   & $-0.210$ small       & No \\
Frustration       & Form vs Ask  & 52 & 403.5 & 0.009   & $-0.414$ medium      & \textbf{Yes} \\
Frustration       & Form vs Chat & 57 & 519.5 & 0.014   & $-0.371$ medium      & \textbf{Yes} \\
Frustration       & Ask vs Chat  & 52 & 629.0 & 0.584   & $+0.087$ negligible  & No \\
AI Reuse ($6_{ai}$) & Form vs Ask  & 49 & 402.5 & 0.034 & $+0.343$ medium      & No \\
AI Reuse ($6_{ai}$) & Form vs Chat & 43 & 216.0 & 0.002 & $+0.543$ large       & \textbf{Yes} \\
AI Reuse ($6_{ai}$) & Ask vs Chat  & 45 & 454.5 & 0.468 & $+0.122$ small       & No \\
\bottomrule
\end{tabular}
\end{table}

\subsection{Mode Preference Ranking}

A chi-square goodness-of-fit test on the rank-1 vote distribution against a
uniform null gives $\chi^2(2, N = 69) = 14.17$, $p < 0.001$, Cramér's
$V = 0.320$ (medium effect). The Friedman test on the full rank scores
($1$--$3$) is non-significant: $\chi^2(2) = 5.07$, $p = 0.079$, $W = 0.037$.

\section{NLP Metric Inferential Tests}
\label{app:inf-nlp}

\subsection{Kruskal--Wallis Omnibus ($n = 113$)}

\begin{table}[htbp]
\centering
\caption{Kruskal--Wallis $H$ tests on the unbalanced full document
set ($n = 113$). Primary omnibus for document-quality inferential
testing.}
\label{tab:app-nlp-kw}
\small
\begin{tabular}{@{}l r r l l@{}}
\toprule
Metric & $H$ & $p$ & Sig ($\alpha = 0.05$) & $\varepsilon^2$ \\
\midrule
BLEU                  & 4.48  & 0.107  & No  & — \\
ROUGE-1               & 5.57  & 0.062  & No  & — \\
ROUGE-2               & 9.56  & 0.008  & Yes & 0.069 (moderate) \\
ROUGE-L               & 7.79  & 0.020  & Yes & — \\
METEOR                & 9.80  & 0.007  & Yes & 0.071 (moderate) \\
BERTScore Precision   & 9.11  & 0.011  & Yes & — \\
BERTScore Recall      & 32.38 & $<0.001$ & Yes & 0.276 (large) \\
BERTScore F1          & 9.75  & 0.008  & Yes & 0.071 (moderate) \\
SBERT (ModernBERT)    & 17.43 & $<0.001$ & Yes & 0.140 (large) \\
\bottomrule
\end{tabular}
\end{table}

\subsection{Pairwise Mann--Whitney $U$ ($n = 113$)}

Pairwise Mann--Whitney $U$ tests on the unbalanced full sample,
serving as the post-hoc family where the Kruskal--Wallis omnibus is
significant. Bonferroni-corrected $\alpha = 0.0167$. Only significant
comparisons are listed; non-significant comparisons are omitted
for brevity.

\begin{table}[htbp]
\centering
\caption{Pairwise Mann--Whitney $U$ post-hoc tests on the
unbalanced full sample of $n = 113$ documents.}
\label{tab:app-nlp-mwu}
\small
\begin{tabular}{@{}l l r r l@{}}
\toprule
Metric & Comparison & $p$ & $r$ & Sig \\
\midrule
BERTScore Recall & Form $>$ Ask  & $<0.001$ & 0.756 (large)  & Yes \\
BERTScore Recall & Form $>$ Chat & $<0.001$ & 0.512 (large)  & Yes \\
BERTScore F1     & Form $>$ Ask  & 0.004    & 0.386 (medium) & Yes \\
BERTScore F1     & Chat $>$ Ask  & 0.011    & 0.340 (medium) & Yes \\
SBERT (ModernBERT) & Chat $>$ Form & $<0.001$ & 0.505 (large)  & Yes \\
SBERT (ModernBERT) & Chat $>$ Ask  & $<0.001$ & 0.453 (large)  & Yes \\
ROUGE-2          & Form $>$ Ask  & 0.006    & 0.369 (medium) & Yes \\
ROUGE-2          & Chat $>$ Ask  & 0.011    & 0.342 (medium) & Yes \\
METEOR           & Form $>$ Ask  & 0.003    & 0.407 (medium) & Yes \\
\bottomrule
\end{tabular}
\end{table}

Ask is the significantly lowest-scoring mode (Ask $<$ Form, Ask
$<$ Chat) on the metrics that reach significance, most clearly on
BERTScore Recall, F1, ROUGE-2, and METEOR; Form and Chat trade
leadership depending on whether the metric emphasizes lexical
recall or document-level semantic alignment.

\section{Confidence--Quality Correlations}
\label{app:inf-corr}

This section reproduces the full pooled and per-mode descriptive
means, Spearman rank correlations, and judge cross-checks used in
the RQ3 triangulation reported in Chapter~\ref{ch:discussion}.

\subsection{Joined Sample and Deduplication}
\label{app:inf-corr-sample}

The triangulation is computed on the joined document-survey
subset: $n = 102$ document-mode rows from $44$ unique participants
(Form $34$, Ask $33$, Chat $35$). The join required a row in both
the NLP/judge file
\texttt{all\_\allowbreak{}scores\_\allowbreak{}combined.csv} (under
\texttt{python/\allowbreak{}metrics/\allowbreak{}output/}) and the
user-study file. Four participants had multiple submissions per
mode in \texttt{document\_results.csv}; the canonical row was
selected by highest \texttt{id} (latest submission). Three
participants present in the documents file had no matching survey
row and were dropped. The resulting $n = 102$ is the source of the
sample size used throughout the RQ3 discussion and differs from
the unmatched join ($n = 106$) reported in earlier drafts. The
underlying tables are available on disk as
\texttt{rq3\_\allowbreak{}confidence\_\allowbreak{}nlp\_\allowbreak{}corr.csv},
\texttt{rq3\_\allowbreak{}per\_\allowbreak{}mode\_\allowbreak{}ranks.csv},
and
\texttt{rq3\_\allowbreak{}judge\_\allowbreak{}vs\_\allowbreak{}confidence\_\allowbreak{}nlp.csv}
under \texttt{python/\allowbreak{}metrics/\allowbreak{}output/} and
\texttt{python/\allowbreak{}judge/\allowbreak{}output/}.

\subsection{Per-Mode Descriptive Means ($n = 102$)}

\begin{table}[htbp]
\centering
\caption{Per-mode means for the confidence, NLP, and judge measures
on the joined RQ3 sample.}
\label{tab:app-corr-means}
\small
\begin{tabular}{@{}l r r r@{}}
\toprule
Measure & Form ($n = 34$) & Ask ($n = 33$) & Chat ($n = 35$) \\
\midrule
Item $5_{ai}$ (legal-basis confidence, 1--5) & $3.38$ & $3.67$ & $3.77$ \\
\texttt{AI\_support} mean (items 1--6, 1--5) & $3.93$ & $3.93$ & $3.88$ \\
BERTScore F1                                 & $0.9059$ & $0.8955$ & $0.9036$ \\
SBERT (ModernBERT)                           & $0.9866$ & $0.9857$ & $0.9895$ \\
Judge legal correctness (1--5)               & $3.59$ & $3.58$ & $3.23$ \\
Judge overall (1--5)                         & $2.38$ & $2.55$ & $2.69$ \\
\bottomrule
\end{tabular}
\end{table}

\subsection{Pooled Spearman Correlations ($n = 102$)}

Primary confidence-quality correlations against the two
reference-based NLP metrics. The pooled correlations are all
essentially zero and none are significant.

\begin{table}[htbp]
\centering
\caption{Pooled Spearman rank correlations between confidence and
reference-based NLP metrics ($n = 102$).}
\label{tab:app-corr-pooled}
\small
\begin{tabular}{@{}l r r r r@{}}
\toprule
& \multicolumn{2}{c}{vs BERTScore F1} & \multicolumn{2}{c}{vs SBERT} \\
\cmidrule(lr){2-3}\cmidrule(lr){4-5}
Confidence measure  & $\rho$ & $p$ & $\rho$ & $p$ \\
\midrule
Item $5_{ai}$              & $-0.013$ & $0.895$ & $+0.004$ & $0.964$ \\
\texttt{AI\_support} mean  & $+0.066$ & $0.510$ & $-0.001$ & $0.995$ \\
\bottomrule
\end{tabular}
\end{table}

\subsection{Per-Mode Spearman Correlations}

Within-mode correlations between item $5_{ai}$ and the two
reference-based metrics. Only Chat shows a positive within-mode
trend, and its correlation against BERTScore F1 sits on the edge
of significance.

\begin{table}[htbp]
\centering
\caption{Per-mode Spearman rank correlations: item $5_{ai}$ versus
reference-based NLP metrics.}
\label{tab:app-corr-permode}
\small
\begin{tabular}{@{}l r r r r r@{}}
\toprule
& & \multicolumn{2}{c}{vs BERTScore F1} & \multicolumn{2}{c}{vs SBERT} \\
\cmidrule(lr){3-4}\cmidrule(lr){5-6}
Mode & $n$ & $\rho$ & $p$ & $\rho$ & $p$ \\
\midrule
Form & 34 & $-0.217$ & $0.22$  & $-0.112$ & $0.53$ \\
Ask  & 33 & $-0.263$ & $0.14$  & $-0.226$ & $0.21$ \\
Chat & 35 & $+0.332$ & $0.051$ & $+0.220$ & $0.20$ \\
\bottomrule
\end{tabular}
\end{table}

\subsection{Judge Cross-Check ($n = 102$)}

The reference-free LLM judge enters as a third lens. Its overall
score correlates strongly with BERTScore F1 and moderately with
SBERT, confirming it tracks the reference-based picture on
holistic quality. The judge's legal-correctness rubric provides
a dimension the reference-based metrics structurally cannot
score; within-mode, Chat is the only mode where confidence and
the judge's verdicts align significantly.

\begin{table}[htbp]
\centering
\caption{Pooled Spearman rank correlations involving the LLM
judge ($n = 102$).}
\label{tab:app-corr-judge-pooled}
\small
\begin{tabular}{@{}l r r@{}}
\toprule
Comparison & $\rho$ & $p$ \\
\midrule
BERTScore F1 vs judge overall            & $+0.566$ & $< 10^{-9}$ \\
SBERT vs judge overall                   & $+0.328$ & $0.0008$ \\
Item $5_{ai}$ vs judge legal correctness & $+0.163$ & $0.10$ \\
Item $5_{ai}$ vs judge overall           & $+0.143$ & $0.15$ \\
\texttt{AI\_support} vs judge overall    & $+0.212$ & $0.03$ \\
\bottomrule
\end{tabular}
\end{table}

\begin{table}[htbp]
\centering
\caption{Chat-only within-mode Spearman rank correlations
involving the LLM judge ($n = 35$). Form and Ask show weak
non-significant trends in the opposite direction.}
\label{tab:app-corr-judge-chat}
\small
\begin{tabular}{@{}l r r@{}}
\toprule
Comparison (Chat only) & $\rho$ & $p$ \\
\midrule
Item $5_{ai}$ vs judge legal correctness & $+0.450$ & $0.007$ \\
Item $5_{ai}$ vs judge overall           & $+0.407$ & $0.015$ \\
\texttt{AI\_support} vs judge overall    & $+0.540$ & $0.0008$ \\
\bottomrule
\end{tabular}
\end{table}

The confidence-correctness inversion identified in
Chapter~\ref{ch:discussion} survives this triangulation at the
mode-rank level (Chat highest on item $5_{ai}$, lowest on
judge legal correctness; Form lowest on item $5_{ai}$, highest on
judge legal correctness) and is absent at the within-participant
level outside Chat.

\section{Timing Inferential Tests}
\label{app:inf-timing}

Inferential tests on the per-mode completion time recorded in the
\texttt{Time elapsed} line emitted by ROPAgen at document submission
(one value per participant--mode pairing). Sample sizes match the NLP
document-level analysis: full $n = 113$ (Form $37$, Ask $37$, Chat
$39$); matched-triple within-subjects subset $n = 26$ participants.
Non-parametric tests are used throughout, consistent with the NLP and
user-study pipelines, as Shapiro--Wilk rejects normality on every
mode (Form $W = 0.780$, $p < 0.001$; Ask $W = 0.912$, $p = 0.030$;
Chat $W = 0.601$, $p < 0.001$). Source:
\texttt{python/\allowbreak{}metrics/\allowbreak{}output/\allowbreak{}statistical\_tests.md}
(\emph{Per-Mode Completion Time} section).

\subsection{Kruskal--Wallis Omnibus ($n = 113$)}

\begin{table}[htbp]
\centering
\caption{Kruskal--Wallis $H$ test on per-mode completion time
across the unbalanced full sample of $n = 113$ documents.}
\label{tab:app-timing-kw}
\small
\begin{tabular}{@{}l r r l l@{}}
\toprule
Measure & $H$ & $p$ & Sig ($\alpha = 0.05$) & $\varepsilon^2$ \\
\midrule
Completion time (s) & 16.455 & $<0.001$ & Yes & 0.147 (large) \\
\bottomrule
\end{tabular}
\end{table}

\subsection{Pairwise Mann--Whitney $U$ ($n = 113$)}

Pairwise Mann--Whitney $U$ tests on the unbalanced full sample,
serving as the post-hoc family where the Kruskal--Wallis omnibus
is significant. Bonferroni-corrected $\alpha = 0.0167$.

\begin{table}[htbp]
\centering
\caption{Pairwise Mann--Whitney $U$ post-hoc tests on per-mode
completion time on the unbalanced full sample of $n = 113$
documents. Effect-size $r$ is positive when the first-named mode
is \emph{slower}.}
\label{tab:app-timing-mwu}
\small
\setlength{\tabcolsep}{4pt}
\begin{tabular}{@{}l r r l r r l@{}}
\toprule
Pair & $U$ & $p$ & $r$ & Median A (s) & Median B (s) & Sig (Bonf) \\
\midrule
Form vs Ask  & 334.0 & $<0.001$ & $+0.512$ large       & 420.0 & 858.0 & \textbf{Yes} \\
Form vs Chat & 420.0 & 0.002    & $+0.418$ medium      & 420.0 & 717.0 & \textbf{Yes} \\
Ask vs Chat  & 793.5 & 0.457    & $-0.100$ negligible  & 858.0 & 717.0 & No \\
\bottomrule
\end{tabular}
\end{table}

Form is significantly faster than both Ask and Chat (large and
medium effects); Ask and Chat are statistically indistinguishable
from one another.
